\documentclass[10pt,sigconf,nonacm]{acmart}

\settopmatter{printacmref=false, printccs=false, printfolios=true}

\usepackage{booktabs} % For formal tables
\usepackage{graphicx} % Required for inserting images

\graphicspath{{figure/}{figures/}}

% Copyright
\setcopyright{none}

%Conference
\acmConference[CSEE E6868]{Embedded Scalable Platforms}{Spring 2026}{New York, NY, USA} 
\acmYear{2026}
\copyrightyear{2026}

\acmPrice{15.00}

\newcommand{\GT}[1]{\textcolor{red}{GT-- #1 --GT}}
\newcommand{\WB}[1]{\textcolor{blue}{WB-- #1 --WB}}

\begin{document}
\title{Intra-Layer Parallelism on Tiled NVDLA Accelerators in ESP}
\subtitle{\normalsize CSEE E6868 - Embedded Scalable Platforms - Spring 2026}

\author{Wenxuan Xu}
\email{xu.wenxuan@columbia.edu} 

\renewcommand{\shortauthors}{W. Xu}


\begin{abstract}
{\small\em
  As deep learning models grow in complexity, running them efficiently on edge devices has become a major challenge. Single-instance accelerators, such as the standard NVIDIA Deep Learning Accelerator (NVDLA), often hit performance walls when handling large workloads. A promising solution is to move away from monolithic designs toward tiled architectures, where multiple smaller accelerators collaborate on executing each layer. This project explores how to implement intra-layer parallelism using a mesh of NVDLA tiles within the Embedded Scalable Platforms (ESP) ecosystem. We plan to instantiate multiple NVDLA tiles connected via a Network-on-Chip (NoC) and—crucially—modify the software stack to split the execution of individual neural network layers across these tiles. By testing different partitioning strategies, such as splitting by output channels or input batches, we aim to evaluate how design-time decisions affect system latency and scalability in tiled architectures.
}
\end{abstract}

\maketitle

\section{Introduction}
\label{sec:intro}
General-purpose processors struggle to keep up with the computational density required by modern Deep Learning (DL) algorithms. This has driven the industry toward specialized hardware accelerators. However, rather than simply building larger, monolithic accelerators, recent research points toward tiled architectures. In this approach, a mesh of smaller, identical accelerator instances cooperate to execute a workload, offering better scalability and yield \cite{survey_accel}. This project investigates a tiled architecture implemented on the ESP research platform \cite{esp_iccad}, where multiple instances of NVDLA \cite{nvdla_web}, an industry-standard open-source accelerator, are integrated to cooperatively execute workloads, enabling a study of system-level performance tradeoffs with increasing tile count.

\subsection{Motivation}
ESP currently supports the integration of a single NVDLA tile \cite{esp_vlsid}, but it does not yet natively support multi-tile execution for a single inference task. Theoretical architectures like NVIDIA's Simba \cite{simba} have demonstrated that a homogeneous mesh of DLAs can achieve significant speedups by exploiting intra-layer parallelism. However, realizing this on an actual FPGA-based SoC involves solving difficult system-level problems. Specifically, we need to map and partition a single layer's computation across multiple independent tiles without the synchronization overhead destroying the performance gains. This project aims to bridge that gap by implementing and evaluating multi-NVDLA execution strategies within ESP.

\section{Specification}
We will leverage the ESP platform's capability to instantiate heterogeneous tiles to build a scalable mesh of NVDLA cores. The target design includes a mesh of NVDLA accelerators (initially a dual-tile setup, potentially scaling to $2 \times 2$ or larger) integrated alongside the mandatory ESP infrastructure, including processor, memory, and I/O tiles.

% --- IMAGE INSERTION POINT ---
\begin{figure}[t]
  \centering
  % Make sure to upload your image and name it 'system_diagram.png'
  \includegraphics[width=\columnwidth]{doc/fig/system_diagram.png}
  \caption{Proposed Architecture: Intra-Layer Parallelism distributing a single Neural Network layer across four NVDLA tiles interconnected via NoC.}
  \label{fig:arch_diagram}
\end{figure}
% -----------------------------

\subsection{Hardware Architecture}
Our baseline relies on the open-source NVDLA RTL \cite{nvdla_web}. We will use the ESP flow to instantiate multiple NVDLA accelerator tiles on the FPGA. The hardware design phase will focus on leveraging the inherent modularity of the ESP architecture to ensure that multiple NVDLA masters can access main memory simultaneously. Specifically, we will utilize ESP's multiplane NoC and instantiate multiple memory tiles to distribute memory traffic and prevent system bus bottlenecks. This configuration is critical for ensuring that the theoretical compute gains are not negated by communication latency.

\subsection{Software and Mapping Strategy}
The biggest challenge in this project is the software stack. The standard NVDLA compiler and runtime are designed for a single hardware target. To enable intra-layer parallelism, we have to intercept the execution flow and partition the workload manually. We intend to explore three main splitting strategies:
\begin{itemize}
    \item \textbf{Output Channel Splitting:} This is our primary strategy. We will assign different subsets of output filters to different tiles, allowing them to compute partial results in parallel.
    \item \textbf{Input Channel Splitting:} Dividing the input feature maps across tiles, which may require more complex synchronization for partial sum accumulation.
    \item \textbf{Batch Splitting:} Processing different samples of a batch in parallel to increase throughput.
\end{itemize}
This will require modifications to the NVDLA Kernel Mode Driver (KMD) and User Mode Driver (UMD) to handle synchronized task submission.

\subsection{Assessment}
We will validate our design using standard benchmarks like ResNet-50 and MobileNet on an FPGA prototype. Key metrics for evaluation include:
\begin{itemize}
    \item \textbf{Functional Correctness:} Ensuring the split execution yields the exact same numerical results as a single-tile run.
    \item \textbf{Scalability:} Quantifying the speedup (or slowdown) as we add more NVDLA tiles.
    \item \textbf{System Efficiency:} Monitoring memory bandwidth usage and NoC congestion to identify bottlenecks.
\end{itemize}

\section{Milestones}
\label{sec:milestones}
We have structured the project to move from a stable single-tile baseline to a multi-tile experimental setup.

\begin{enumerate}
  \item \textbf{Baseline Setup \& NVDLA Study (by Feb. 23)}
  \begin{itemize}
      \item Review NVDLA documentation \cite{nvdla_web} and the Simba architecture paper \cite{simba}.
      \item Replicate the single-NVDLA tutorial on FPGA to verify the toolchain.
  \end{itemize}
  
  \item \textbf{Multi-NVDLA Instantiation (by Mar. 11)}
  \begin{itemize}
      \item Configure an ESP SoC with two NVDLA tiles.
      \item Verify the hardware stability by running two separate networks simultaneously (independent execution).
  \end{itemize}
  
  \item \textbf{Mid-term Presentation and Report ($\sim$ Mar. 11)}
  
  \item \textbf{Intra-Layer Partitioning Implementation (by Mar. 30)}
  \begin{itemize}
      \item Implement the software logic to split a single layer across two tiles.
      \item Focus first on "Output Channel" splitting as the proof-of-concept.
      \item Debug synchronization and interrupt handling between the two tiles.
  \end{itemize}
  
  \item \textbf{Design Space Exploration (by Apr. 20)}
  \begin{itemize}
      \item If resources allow, scale the design to a larger mesh (e.g., $2 \times 2$).
      \item Compare performance differences between Input vs. Output channel splitting.
      \item Collect bandwidth and latency data using ESP's performance counters.
  \end{itemize}
  
  \item \textbf{Final Evaluation \& Analysis (by May 4)}
  \begin{itemize}
      \item Conduct ablation studies to isolate the impact of NoC configuration.
      \item Finalize performance plots and analysis.
  \end{itemize}
  
  \item \textbf{Final presentation and report ($\sim$ May 8 - 15)}
\end{enumerate}

\section{Critical Aspects}
\textbf{Software Stack Complexity:} The NVDLA software stack is mature but notoriously complex. Modifying the runtime to handle fragmented layer execution without rewriting the entire compiler is a significant risk. We plan to mitigate this by manipulating the workload descriptor lists at the driver level rather than touching the compiler front-end.

\textbf{Memory Bandwidth:} NVDLA is bandwidth-hungry. There is a risk that scaling to multiple tiles will saturate the on-chip NoC bandwidth and/or the shared off-chip memory interface of our FPGA prototype, masking any compute gains. We will need to use ESP's monitoring tools early on to detect if bandwidth becomes the bottleneck.

\begin{thebibliography}{1}

\bibitem{simba}
Y. S. Shao, J. Clemons, R. Venkatesan, B. Zimmer, M. Fojtik, N. Jiang, B. Keller, A. Klinefelter, N. Pinckney, P. Raina, S. G. Tell, Y. Zhang, W. J. Dally, J. Emer, C. T. Gray, B. Khailany, and S. W. Keckler, "Simba: Scaling Deep-Learning Inference with Multi-Chip-Module-Based Architecture," in \emph{Proceedings of the 52nd Annual IEEE/ACM International Symposium on Microarchitecture (MICRO)}, 2019.

\bibitem{nvdla_web}
NVIDIA Corporation, "NVDLA Open Source Project," http://nvdla.org/, Accessed: 2026.

\bibitem{esp_iccad}
P. Mantovani, D. Giri, G. Di Guglielmo, L. Piccolboni, J. Zuckerman, E. G. Cota, M. Petracca, C. Pilato, and L. P. Carloni, "Agile SoC Development with Open ESP," in \emph{IEEE/ACM International Conference on Computer Aided Design (ICCAD)}, 2020.

\bibitem{esp_vlsid}
L. P. Carloni, "ESP: The Open-Source SoC Platform," in \emph{International Conference on VLSI Design and International Conference on Embedded Systems (VLSID)}, 2020.

\bibitem{survey_accel}
A. Reuther, P. Michaleas, M. Jones, V. Gadepally, S. Samsi, and J. Kepner, "Survey of Machine Learning Accelerators," in \emph{IEEE High Performance Extreme Computing Conference (HPEC)}, 2020.

\end{thebibliography}

\end{document}